% Options for packages loaded elsewhere
% Options for packages loaded elsewhere
\PassOptionsToPackage{unicode}{hyperref}
\PassOptionsToPackage{hyphens}{url}
\PassOptionsToPackage{dvipsnames,svgnames,x11names}{xcolor}
%
\documentclass[
  spanish,
  12pt,
  letterpaper,
]{article}
\usepackage{xcolor}
\usepackage[top=2.5cm,bottom=2.5cm,left=2.5cm,right=2.5cm]{geometry}
\usepackage{amsmath,amssymb}
\setcounter{secnumdepth}{5}
\usepackage{iftex}
\ifPDFTeX
  \usepackage[T1]{fontenc}
  \usepackage[utf8]{inputenc}
  \usepackage{textcomp} % provide euro and other symbols
\else % if luatex or xetex
  \usepackage{unicode-math} % this also loads fontspec
  \defaultfontfeatures{Scale=MatchLowercase}
  \defaultfontfeatures[\rmfamily]{Ligatures=TeX,Scale=1}
\fi
\usepackage{lmodern}
\ifPDFTeX\else
  % xetex/luatex font selection
\fi
% Use upquote if available, for straight quotes in verbatim environments
\IfFileExists{upquote.sty}{\usepackage{upquote}}{}
\IfFileExists{microtype.sty}{% use microtype if available
  \usepackage[]{microtype}
  \UseMicrotypeSet[protrusion]{basicmath} % disable protrusion for tt fonts
}{}
\makeatletter
\@ifundefined{KOMAClassName}{% if non-KOMA class
  \IfFileExists{parskip.sty}{%
    \usepackage{parskip}
  }{% else
    \setlength{\parindent}{0pt}
    \setlength{\parskip}{6pt plus 2pt minus 1pt}}
}{% if KOMA class
  \KOMAoptions{parskip=half}}
\makeatother
% Make \paragraph and \subparagraph free-standing
\makeatletter
\ifx\paragraph\undefined\else
  \let\oldparagraph\paragraph
  \renewcommand{\paragraph}{
    \@ifstar
      \xxxParagraphStar
      \xxxParagraphNoStar
  }
  \newcommand{\xxxParagraphStar}[1]{\oldparagraph*{#1}\mbox{}}
  \newcommand{\xxxParagraphNoStar}[1]{\oldparagraph{#1}\mbox{}}
\fi
\ifx\subparagraph\undefined\else
  \let\oldsubparagraph\subparagraph
  \renewcommand{\subparagraph}{
    \@ifstar
      \xxxSubParagraphStar
      \xxxSubParagraphNoStar
  }
  \newcommand{\xxxSubParagraphStar}[1]{\oldsubparagraph*{#1}\mbox{}}
  \newcommand{\xxxSubParagraphNoStar}[1]{\oldsubparagraph{#1}\mbox{}}
\fi
\makeatother


\usepackage{longtable,booktabs,array}
\newcounter{none} % for unnumbered tables
\usepackage{calc} % for calculating minipage widths
% Correct order of tables after \paragraph or \subparagraph
\usepackage{etoolbox}
\makeatletter
\patchcmd\longtable{\par}{\if@noskipsec\mbox{}\fi\par}{}{}
\makeatother
% Allow footnotes in longtable head/foot
\IfFileExists{footnotehyper.sty}{\usepackage{footnotehyper}}{\usepackage{footnote}}
\makesavenoteenv{longtable}
\usepackage{graphicx}
\makeatletter
\newsavebox\pandoc@box
\newcommand*\pandocbounded[1]{% scales image to fit in text height/width
  \sbox\pandoc@box{#1}%
  \Gscale@div\@tempa{\textheight}{\dimexpr\ht\pandoc@box+\dp\pandoc@box\relax}%
  \Gscale@div\@tempb{\linewidth}{\wd\pandoc@box}%
  \ifdim\@tempb\p@<\@tempa\p@\let\@tempa\@tempb\fi% select the smaller of both
  \ifdim\@tempa\p@<\p@\scalebox{\@tempa}{\usebox\pandoc@box}%
  \else\usebox{\pandoc@box}%
  \fi%
}
% Set default figure placement to htbp
\def\fps@figure{htbp}
\makeatother



\ifLuaTeX
\usepackage[bidi=basic,shorthands=off]{babel}
\else
\usepackage[bidi=default,shorthands=off]{babel}
\fi
\ifLuaTeX
  \usepackage{selnolig} % disable illegal ligatures
\fi


\setlength{\emergencystretch}{3em} % prevent overfull lines

\providecommand{\tightlist}{%
  \setlength{\itemsep}{0pt}\setlength{\parskip}{0pt}}



 


\makeatletter
\@ifpackageloaded{caption}{}{\usepackage{caption}}
\AtBeginDocument{%
\ifdefined\contentsname
  \renewcommand*\contentsname{Tabla de contenidos}
\else
  \newcommand\contentsname{Tabla de contenidos}
\fi
\ifdefined\listfigurename
  \renewcommand*\listfigurename{Listado de Figuras}
\else
  \newcommand\listfigurename{Listado de Figuras}
\fi
\ifdefined\listtablename
  \renewcommand*\listtablename{Listado de Tablas}
\else
  \newcommand\listtablename{Listado de Tablas}
\fi
\ifdefined\figurename
  \renewcommand*\figurename{Figura}
\else
  \newcommand\figurename{Figura}
\fi
\ifdefined\tablename
  \renewcommand*\tablename{Tabla}
\else
  \newcommand\tablename{Tabla}
\fi
}
\@ifpackageloaded{float}{}{\usepackage{float}}
\floatstyle{ruled}
\@ifundefined{c@chapter}{\newfloat{codelisting}{h}{lop}}{\newfloat{codelisting}{h}{lop}[chapter]}
\floatname{codelisting}{Listado}
\newcommand*\listoflistings{\listof{codelisting}{Listado de Listados}}
\makeatother
\makeatletter
\makeatother
\makeatletter
\@ifpackageloaded{caption}{}{\usepackage{caption}}
\@ifpackageloaded{subcaption}{}{\usepackage{subcaption}}
\makeatother
\usepackage{bookmark}
\IfFileExists{xurl.sty}{\usepackage{xurl}}{} % add URL line breaks if available
\urlstyle{same}
\hypersetup{
  pdflang={es},
  colorlinks=true,
  linkcolor={blue},
  filecolor={Maroon},
  citecolor={Blue},
  urlcolor={Blue},
  pdfcreator={LaTeX via pandoc}}


\author{}
\date{}
\begin{document}


\begin{flushleft}

{\small
\textbf{Universidad de Costa Rica}\\
Escuela de Matemática\\
CA-0204 — Herramientas de Ciencia de Datos I\\
II ciclo, 2026
}

\vspace{0.3cm}

{\Large\textbf{StoryTrack}}\\
\textbf{Avance 3 — Informe de resultados}

\vspace{0.25cm}

{\small
Dylan Madrigal Jiménez \quad María José Angulo Salas\\
Gabriel Guevara Cornejo \quad Harol Arias Sánchez
}

\vspace{0.15cm}

{\small 6 de noviembre de 2026}

\end{flushleft}

\noindent

\rule{\linewidth}{0.4pt}

\subsection{Introducción}\label{introducciuxf3n}

Las series de televisión pueden contar con numerosas temporadas y
episodios, lo que dificulta recordar todos los acontecimientos
importantes de la historia. Esto ocurre especialmente cuando una persona
deja de ver una serie durante un timepo y, al retomarla, no recuerda con
claridad lo sucedido en episodios anteriores.

Ante esta situación surge \textbf{Story Track}, un proyecto que busca
desarrollar una herramienta que permite a los usuarios llevar un
seguimiento de las series que ven, consultar su progreso y acceder a
resúmenes de los episodios anteriores. Su finalidad es facilitar que las
personas puedan retomar una serie sin necesidad de volver a ver
capítulos completos para recordar los acontecimientos más relevantes.

Para desarrollar el proyecto se utiliza el lenguaje de programación R
junto con herramientas de inteligencia artificial. Mediante R, se extrae
y organiza la información de los archivos de subtítulos en formato .srt,
mientras que la inteligencia artificial se utiliza para analizar y
resumir el contenido de los diálogos. Esto permitirá generar resúmenes
organizados por episodio, que el usuario podrá consultar según su avance
en la serie.

De esta manera, StoryTrack busca ofrecer una solución práctica para dar
seguimiento a las series de televisión, combinando el procesamiento de
datos y la inteligencia artificial para facilitar el acceso a
información sobre los episodios y sus acontecimientos principales.

\subsection{Objetivo general}\label{objetivo-general}

Desarrollar una herramienta que facilite el segumiento de series de
televisión, permitiendo al usuario consultar su progreso y recordar los
acontecimientos más importantes de la historia mediante resúmenes
generados con inteligencia artificial.

\subsection{Metodología}\label{metodologuxeda}

\section{Obtención y procesamiento de los
subtítulos}\label{obtenciuxf3n-y-procesamiento-de-los-subtuxedtulos}

Para el desarrollo de StoryTrack se utilizaron únicamente los subtítulos
de la serie \emph{The Big Bang Theory}, con el propósito de trabajar con
una base de datos más sencilla y delimitar el alcance del proyecto. Esto
permitió concentrarse en las funciones principales de la herramienta.
Los archivos utilizados se encuentran en formato .srt, el cual contiene
los diálogos de los episodios junto con sus respectivos tiempos de
aparición.

El procesamiento se realizó mediante el lenguaje de programación R. Para
ello, se creó la función leer.subtitulos (), encargada de leer los
archivos, identificar los bloques de subtítulos y extraer la información
necesaria para construir una base de datos estructurada.

A partir de los bloques de subtítulos se extrajeron nueve variables:
nombre de la serie, temporada, episodio, número de bloque, tiempo de
inicio, tiempo de finalización, personaje, texto y nombre del archivo
original. La variable personaje se incorporó posteriormente para
identificar a quién corresponde cada diálogo, cuando esta información se
encuentra disponible en los subtítulos.

\section{Filtrado de la historia}\label{filtrado-de-la-historia}

Posteriormente, se desarrolló la función filtrar.historia(), con el
propósito de limitar la información disponible según el progreso del
usuario en la serie y evitar que se revelen acontecimientos que todavía
no ha visto.

Esta función recibe la base de datos de subtítulos y la información
correspondiente a la serie, temporada y tiempo seleccionado por el
usuario. Para realizar el filtrado, los tiempos de los subtítulos se
convierten a segundos, lo que permite compararlos con el límite
establecido.

De esta manera, se conservan los subtítulos de las temporadas y
episodios anteriores al punto seleccionado, así como aquellos del
episodio actual cuyo tiempo de finalización no supera el minuto
indicado. Este procedimiento constituye la base del sistema denominado
\textbf{Spoiler Guard}, encargado de restringir la información utilizada
posterioremente por la inteligencia artificial.

\section{Creación del contexto}\label{creaciuxf3n-del-contexto}

Una vez filtrada la información, se desarrolló la función
crear.contexto(), encargada de organizar los subtítulos seleccionados en
un formato de texto que pueda utilizarse como entrada para la
inteligencia artificial.

Para ello, se combinan los datos de cada regsitro, incluyendo la
temporada, el episodio, los tiempos de inicio y finalización, el
personaje identificado y el diálogo correspondiente.

Posteriormente, los registros se unen en un único texto, se unen en un
único texto, conservando la información necesaria para identificar a qué
episodio y momento pertenece cada intervención. Así, se construye el
contexto que se enviará a la inteligencia artificial, utilizando
únicamente los subtítulos permitidos por el filtrado anterior.

\section{Desarrollo de la aplicación en
Shiny}\label{desarrollo-de-la-aplicaciuxf3n-en-shiny}

Para integrar las funciones anteriores, se desarrolló una aplicacion
interactiva mediante la libreria shiny de R, cuya estructura principal
se encuentra en el archivo app.R.

Dentro de este archivo se incorporó la lectura y consolidación de los
subtítulos procesados, así como un catálogo de series y una interfaz que
permite seleccionar temporadas y episodios. Aunque el catálogo incluye
otras series como ejemplos, el procesamiento de subtítulos se desarrolló
utilizando \emph{The Big Bang Theory}. La aplicación permite al usuario
registrar su progreso.

Se integraron las funciones filtrar.historia () y crear.contexto(), de
manera que la información utilizada para generar resúmenes o responder
presguntas se ajuste al progreso registrado por el usuario. La interfaz
también incorpora las opciones para solicitar resúmenes y realizar
preguntas relacionadas con los personajes y acontecimientos de la serie.

\section{Integración con ChatGPT}\label{integraciuxf3n-con-chatgpt}

Finalmente, se desarrolló el archivo chatgpt\_local.R, encargado de
gestionar la conexión entre StoryTrack y ChatGPT. Este código implementa
un procedimiento de autenticación para que el usuario pueda autorizar el
acceso desde la aplicación.

Una vez establecida la conexión, la aplicación está programada para
enviar a ChatGPT el contexto previamente filtrado, junto con
instrucciones específicas para generar resúmenes o responder preguntas
sobre la serie.

Las instrucciones establecen que las respuestas deben elaborarse
utilizando únicamente la información proporcionada en los subtítulos,
sin incorporar acontecimientos posteriores al progreso del usuario.
Asimismo, se indica que, cuando la información disponible no permita
responder una pregunta, debe señalarse esta limitación.

De esta forma, se busca combinar el seguimiento de episodios con la
generación de resúmenes y respuestas mediante inteligencia artificial,
utilizando SpoilerGuard para limitar el contenido enviado según el
avance del usuario.

\section{Resultados}\label{resultados}

{[}PENDIENTE: Análisis, gráficos y tablas.{]}

\section{Bibliografía}\label{bibliografuxeda}

{[}PENDIENTE: Referencias APA 7.{]}

\section{Anexos}\label{anexos}

{[}PENDIENTE: Material complementario y enlace al repositorio.{]}




\end{document}
